% 原试卷 PDF 第46页。
\begin{question}[index=11]
$\displaystyle\lim_{x\to0}\left(\frac{2+x}{\sin x}-\frac2x\right)=$\fillin。
\end{question}
\QuestionBreak

\begin{question}[index=12]
设可微函数 $f(x)$ 及其反函数 $g(x)$ 满足
\[
\int_1^{f(x)}g(t)\dd t=\frac13\left(x^{\frac32}-8\right)\quad(x>0)\text{，}
\]
则 $f(x)=$\fillin。
\end{question}
\QuestionBreak

\begin{question}[index=13]
$\displaystyle\sum_{n=1}^{\infty}(-1)^n\frac{n}{2^n}(x-1)^{2n}$ 的收敛域为\fillin。
\end{question}
\QuestionBreak

\begin{question}[index=14]
$\displaystyle\int_{-2}^2\dd x\int_{-\sqrt{4-x^2}}^{\sqrt{4-x^2}}\dd y
\int_{(x^2+y^2)^2}^0 x^2\dd z=$\fillin。
\end{question}
\QuestionBreak

\begin{question}[index=15]
二次型 $f(x_1,x_2,x_3)=(x_1-x_2+x_3)^2+(x_2-ax_3)^2+(ax_3+x_1)^2$ 的秩为 $2$，则 $a=$\fillin。
\end{question}
\QuestionBreak

\begin{question}[index=16]
设 $A,B,C$ 为随机事件，$A$ 与 $B$ 相互独立，且 $P(A)=P(B)=P(C)=\dfrac13$，$P(B\cup C)=\dfrac23$，则 $P(AB\cup C)=$\fillin。
\end{question}
